\chapter{Girsanov et changements de numéraire}\label{ch:qm:girsanov-and-changes-of-numeraire}

Un système de gestion des risques évalue un caplet à cinq ans en simulant le taux court semaine après semaine,
en l'intégrant dans un facteur d'actualisation et en moyennant sur les trajectoires~: 260 nombres aléatoires par trajectoire et une
erreur type de 0.05~pb après 20\,000 trajectoires. Un collègue évalue le même caplet en changeant de probabilité pour celle
attachée à l'obligation zéro-coupon échéant à la date de fixing~:
sous cette probabilité, le taux court à cette date est une unique gaussienne de moyenne connue, de sorte que chaque trajectoire coûte
un seul nombre aléatoire, et le même raisonnement donne une formule fermée qui prend une microseconde. La
réponses concordent à l'intérieur de leurs erreurs types~: 3.26~points de base de notionnel. Rien n'a changé sur le
marché~; seule l'unité dans laquelle on compte les valeurs a changé. Ce chapitre démontre le
théorème de Girsanov, qui dit ce que fait un \omterm{def:qm:probability-at-speed:change}{changement de probabilité} à un \omterm{def:qm:brownian-motion:bm}{mouvement brownien}, définit les \omterm{def:qm:girsanov-and-changes-of-numeraire:numeraire}{numéraires} et
les probabilités \omterm{def:qm:probability-at-speed:martingale}{martingales} qui leur sont attachées, et établit les \omterm{def:qm:girsanov-and-changes-of-numeraire:forward}{probabilités forward} et annuité que
calculent tous les chapitres de la série consacrés aux taux et aux options.

\section{Théorème de Girsanov}

\begin{definition}[Condition de Novikov]\label{def:qm:girsanov-and-changes-of-numeraire:novikov}
Un \omterm{def:qm:probability-at-speed:filtration}{processus adapté} $\gamma$ vérifie la \emph{condition de Novikov}\index{condition de Novikov} (Novikov's condition) sur $[0,
T]$ si $\E[\exp(\tfrac12\int_0^T\gamma_s^2\,ds)] < \infty$.
\end{definition}

C'est une condition suffisante pour que la \omterm{def:qm:ito-calculus:local}{martingale locale} positive $Z_t = \mathcal E(\int\gamma\,dW)_t =
\exp(\int_0^t\gamma_s\,dW_s - \tfrac12\int_0^t\gamma_s^2\,ds)$ de la \cref{def:qm:ito-calculus:exp} soit une
véritable \omterm{def:qm:probability-at-speed:martingale}{martingale}, avec $\E[Z_T] = 1$, de sorte qu'elle puisse servir de \omterm{def:qm:probability-at-speed:change}{processus densité}
(\cref{def:qm:probability-at-speed:change}). Un $\gamma$ borné vérifie toujours cette condition.

\begin{theorem}[Girsanov]\label{thm:qm:girsanov-and-changes-of-numeraire:girsanov}
Soit $\gamma$ vérifiant la \omterm{def:qm:girsanov-and-changes-of-numeraire:novikov}{condition de Novikov}, et posons $\mathbb Q$ sur $\mathcal F_T$ par $d\mathbb Q/d\P = Z_T$. Alors
\[
W^{\mathbb Q}_t = W_t - \int_0^t\gamma_s\,ds
\]
est un \omterm{def:qm:brownian-motion:bm}{mouvement brownien} sous $\mathbb Q$. Réciproquement, toute $\mathbb Q \sim \P$ sur la \omterm{def:qm:probability-at-speed:filtration}{filtration} de $W$ admet une densité
de cette forme.
\end{theorem}

\begin{proof}
Par le \cref{thm:qm:ito-calculus:levy}, il suffit que $W^{\mathbb Q}$ soit une $\mathbb Q$-martingale locale continue avec
$[W^{\mathbb Q}]_t = t$~; la \omterm{def:qm:brownian-motion:qv}{variation quadratique} est celle de $W$, puisque la dérive est à variation finie.
Par la \cref{prop:qm:probability-at-speed:bayes}, $W^{\mathbb Q}$ est une $\mathbb Q$-martingale locale si $ZW^{\mathbb Q}$ est une
$\P$-martingale locale, et par la règle du produit, avec $dZ = Z\gamma\,dW$,
$d(ZW^{\mathbb Q}) = Z\,dW^{\mathbb Q} + W^{\mathbb Q}\,dZ + d[Z, W^{\mathbb Q}] = Z\,dW - Z\gamma\,dt + W^{\mathbb Q}\,dZ + Z\gamma\,dt$, une
intégrale stochastique. La réciproque est le théorème de représentation des \omterm{def:qm:probability-at-speed:martingale}{martingales} appliqué au
\omterm{def:qm:probability-at-speed:change}{processus densité} (\cref{thm:qm:ito-calculus:representation}).
\end{proof}

Un \omterm{def:qm:probability-at-speed:change}{changement de probabilité} ne peut pas modifier la volatilité, ce qui se voit dans la \omterm{def:qm:brownian-motion:qv}{variation quadratique} d'une
trajectoire isolée~; il modifie en revanche la dérive, qui n'est pas invariante. Avec $\gamma$ constant, le théorème est la
formule de Cameron--Martin~: pondérer les trajectoires par $e^{\gamma W_T - \gamma^2T/2}$ transforme un \omterm{def:qm:brownian-motion:bm}{mouvement brownien} en un
\omterm{def:qm:brownian-motion:bm}{mouvement brownien} de dérive $\gamma$ (\cref{fig:qm:girsanov-and-changes-of-numeraire:girsanov}). Lu à rebours, il supprime une
dérive, ce qui permet d'établir la probabilité de premier passage avec dérive du chapitre~2
(l'\cref{exo:qm:girsanov-and-changes-of-numeraire:7}).

\begin{omfigure}
\begin{tikzpicture}
\begin{axis}[width=11cm,height=4.8cm,
  xlabel={$x$},ylabel={densité de $W_1$},
  tick label style={font=\small},label style={font=\small},
  xmin=-4,xmax=5,ymin=0,ymax=0.45,grid=major,grid style={gray!25},
  legend columns=2,legend style={font=\footnotesize,at={(0.5,-0.3)},anchor=north,draw=none,fill=none,/tikz/every even column/.append style={column sep=8pt}},
  table/col sep=comma]
\addplot[omDef,thick,const plot mark mid,mark=none] table[x=x,y=p]{figdata/methods/05-girsanov-and-changes-of-numeraire/girsanov.csv};
\addplot[omThm,thick,const plot mark mid] table[x=x,y=q]{figdata/methods/05-girsanov-and-changes-of-numeraire/girsanov.csv};
\addplot[omDef,dashed,domain=-4:5,samples=80,forget plot] {exp(-x^2/2)/sqrt(2*pi)};
\addplot[omThm,dashed,domain=-4:5,samples=80,forget plot] {exp(-(x-1)^2/2)/sqrt(2*pi)};
\legend{{sous $\P$~: $\mathcal N(0,1)$},{repondérée par $e^{W_1 - 1/2}$~: $\mathcal N(1,1)$}}
\end{axis}
\end{tikzpicture}

\omcaption{Les mêmes 200\,000 tirages de $W_1$, comptés simplement (en bleu) et pondérés par la densité
$Z_1 = e^{\gamma W_1 - \gamma^2/2}$ avec $\gamma = 1$ (en rouge)~: l'histogramme pondéré est la densité de $\mathcal N(1, 1)$
(tirets), et la forme reste inchangée. Données~: tutoriel du chapitre, avec graine fixée.}\label{fig:qm:girsanov-and-changes-of-numeraire:girsanov}
\end{omfigure}

\section{Numéraires et leurs probabilités martingales}

\begin{definition}[Numéraire, compte monétaire]\label{def:qm:girsanov-and-changes-of-numeraire:numeraire}
Un \emph{numéraire}\index{numéraire} $\mathcal N$ est le processus de prix d'un actif négocié (ou d'un
portefeuille autofinancé) qui est strictement positif à tout instant~; les valeurs exprimées en unités de celui-ci sont
$V_t/\mathcal N_t$. Le \emph{compte monétaire}\index{compte monétaire} (money-market account) est le numéraire $B_t =
\exp\int_0^tr_s\,ds$ qui se capitalise au jour le jour au taux court $r_t$.
\end{definition}

\begin{definition}[Probabilité martingale équivalente, probabilité risque-neutre]\label{def:qm:girsanov-and-changes-of-numeraire:emm}
Une \emph{probabilité martingale équivalente}\index{probabilité martingale équivalente} (equivalent \omterm{def:qm:probability-at-speed:martingale}{martingale} measure) pour un \omterm{def:qm:girsanov-and-changes-of-numeraire:numeraire}{numéraire} $\mathcal N$
est une probabilité $\mathbb Q^{\mathcal N}$ équivalente à $\P$ sous laquelle $V_t/\mathcal N_t$ est une \omterm{def:qm:probability-at-speed:martingale}{martingale} pour
le prix $V$ de tout actif négocié. La \emph{probabilité risque-neutre}\index{probabilité risque-neutre} (risk-neutral measure)
$\mathbb Q$ est la probabilité \omterm{def:qm:probability-at-speed:martingale}{martingale} équivalente pour le \omterm{def:qm:girsanov-and-changes-of-numeraire:numeraire}{compte monétaire}.
\end{definition}

Sous une telle probabilité, un actif versant $V_T$ en $T$ vaut
\[
V_t = \mathcal N_t\,\E^{\mathbb Q^{\mathcal N}}_t\Bigl[\frac{V_T}{\mathcal N_T}\Bigr], \qquad\text{et avec } \mathcal N = B:\quad V_t =
\E^{\mathbb Q}_t\bigl[e^{-\int_t^Tr_s\,ds}V_T\bigr].
\]
Pourquoi une telle probabilité devrait exister relève d'une affirmation économique~: dans un marché sans opportunité d'arbitrage
il en existe au moins une, et dans un marché complet exactement une. C'est le théorème fondamental de
l'évaluation d'actifs, et c'est le livre~5, chapitre~1, qui le démontre~; ici les probabilités sont admises
et utilisées comme des outils. Sous $\mathbb Q$, une action qui ne verse pas de dividende a pour dérive $r_t$~: $dS/S = r_t\,dt +
\sigma\,dW^{\mathbb Q}$, quel que soit son rendement espéré dans le monde réel, car Girsanov avec $\gamma = -(\mu -
r)/\sigma$ élimine l'excédent.

\section{Changement de numéraire}

\begin{definition}[Changement de numéraire]\label{def:qm:girsanov-and-changes-of-numeraire:changenum}
Un \emph{changement de numéraire}\index{changement de numéraire} (change of numeraire) remplace un \omterm{def:qm:girsanov-and-changes-of-numeraire:numeraire}{numéraire} et sa \omterm{def:qm:girsanov-and-changes-of-numeraire:emm}{probabilité martingale équivalente} par un autre \omterm{def:qm:girsanov-and-changes-of-numeraire:numeraire}{numéraire} et la probabilité sous laquelle les prix divisés par celui-ci sont des
\omterm{def:qm:probability-at-speed:martingale}{martingales}~; les prix ne changent pas, seuls changent l'unité et la probabilité dans lesquelles ils sont calculés.
\end{definition}

\begin{theorem}[Changement de numéraire]\label{thm:qm:girsanov-and-changes-of-numeraire:change}
Si $\mathbb Q^{\mathcal N}$ est une \omterm{def:qm:girsanov-and-changes-of-numeraire:emm}{probabilité martingale équivalente} pour $\mathcal N$ et si $\mathcal M$ est un autre
\omterm{def:qm:girsanov-and-changes-of-numeraire:numeraire}{numéraire}, la probabilité définie par
\[
\frac{d\mathbb Q^{\mathcal M}}{d\mathbb Q^{\mathcal N}}\bigg|_{\mathcal F_t} = \frac{\mathcal M_t/\mathcal M_0}{\mathcal N_t/\mathcal N_0}
\]
est une \omterm{def:qm:girsanov-and-changes-of-numeraire:emm}{probabilité martingale équivalente} pour $\mathcal M$. Si $d\mathcal M/\mathcal M$ et $d\mathcal N/\mathcal N$ ont des
vecteurs de volatilité $\sigma_{\mathcal M}$, $\sigma_{\mathcal N}$ par rapport à un \omterm{def:qm:brownian-motion:bm}{mouvement brownien} $W^{\mathcal N}$ de $\mathbb
Q^{\mathcal N}$, alors $dW^{\mathcal M} = dW^{\mathcal N} - (\sigma_{\mathcal M} - \sigma_{\mathcal N})\,dt$.
\end{theorem}

\begin{proof}
Le \omterm{def:qm:probability-at-speed:change}{processus densité} $Z_t = (\mathcal M_t/\mathcal N_t)(\mathcal N_0/\mathcal M_0)$ est une $\mathbb Q^{\mathcal N}$-martingale positive
car $\mathcal M$ est un actif négocié. Pour un actif négocié $V$, $Z_t\cdot V_t/\mathcal M_t = (\mathcal N_0/\mathcal M_0)\,V_t/\mathcal N_t$ est une
$\mathbb Q^{\mathcal N}$-martingale, donc $V/\mathcal M$ est une $\mathbb Q^{\mathcal M}$-martingale d'après
la \cref{prop:qm:probability-at-speed:bayes}. Par la formule d'Itô, le rapport $\mathcal M/\mathcal N$ a pour volatilité
$\sigma_{\mathcal M} - \sigma_{\mathcal N}$, donc $Z = \mathcal E(\int(\sigma_{\mathcal M} - \sigma_{\mathcal N})\,dW^{\mathcal N})$ et Girsanov
donne la dérive.
\end{proof}

Geman, El Karoui et Rochet (1995) ont fait de cet outil un instrument de travail courant. L'art consiste à choisir le
\omterm{def:qm:girsanov-and-changes-of-numeraire:numeraire}{numéraire} pour que le rapport $V_T/\mathcal N_T$ soit simple. La \Cref{fig:qm:girsanov-and-changes-of-numeraire:map}
recense les choix retenus dans la série~; l'option de Margrabe d'échange d'un actif contre un autre, valorisée
avec le second actif comme \omterm{def:qm:girsanov-and-changes-of-numeraire:numeraire}{numéraire}, fait l'objet de l'\cref{exo:qm:girsanov-and-changes-of-numeraire:6}.

\begin{omfigure}
\begin{tikzpicture}[font=\small]
\matrix[matrix of nodes,nodes={align=left,anchor=west,inner sep=3pt},
  column 1/.style={nodes={text width=3.6cm}},column 2/.style={nodes={text width=3.4cm}},
  column 3/.style={nodes={text width=5.6cm}},row sep=4pt,column sep=6pt] (m) {
\textbf{numéraire} & \textbf{probabilité} & \textbf{martingales (prix$\,/\,$numéraire)} \\
compte du marché monétaire $B_t$ & risque-neutre $\mathbb Q$ & tout prix actualisé $V_t/B_t$ \\
obligation zéro-coupon $P(t,T)$ & $T$-forward $\mathbb Q^T$ & prix forward $V_t/P(t,T)$~; un taux simple $F(t;T{-}\delta,T)$ versé en $T$ \\
annuité $A_t$ & probabilité annuité $\mathbb Q^A$ & le taux de swap pair $S_t$ du swap de dates fixes $T_i$ \\
une action $S_t$ & probabilité action $\mathbb Q^S$ & $V_t/S_t$~: le ratio $S^2_t/S^1_t$ de l'option d'échange \\
};
\draw[omDef,thick] (m-1-1.north west) -- (m-1-3.north east);
\draw[omDef] (m-1-1.south west) -- (m-1-3.south east);
\draw[omDef,thick] (m-5-1.south west) -- (m-5-3.south east);
\end{tikzpicture}

\omcaption{\omterm{def:qm:girsanov-and-changes-of-numeraire:numeraire}{Numéraires} utilisés dans la série, leurs \omterm{def:qm:girsanov-and-changes-of-numeraire:emm}{probabilités martingales équivalentes}, et ce qui devient une
\omterm{def:qm:probability-at-speed:martingale}{martingale} sous chacune d'elles~: valoriser un produit revient à choisir le \omterm{def:qm:girsanov-and-changes-of-numeraire:numeraire}{numéraire} qui rend son rapport
le plus simple.}\label{fig:qm:girsanov-and-changes-of-numeraire:map}
\end{omfigure}

\section{La probabilité forward}

\begin{definition}[Probabilité forward]\label{def:qm:girsanov-and-changes-of-numeraire:forward}
La \emph{probabilité forward}\index{probabilité forward} $\mathbb Q^T$ d'échéance $T$ est la \omterm{def:qm:girsanov-and-changes-of-numeraire:emm}{probabilité martingale
équivalente} pour le \omterm{def:qm:girsanov-and-changes-of-numeraire:numeraire}{numéraire} $P(t, T)$, l'obligation zéro-coupon arrivant à échéance en $T$.
\end{definition}

\begin{proposition}[Actualisation dans l'espérance et hors de celle-ci]\label{prop:qm:girsanov-and-changes-of-numeraire:forward}
Pour un gain $X$ en $T$~: $\E^{\mathbb Q}_t[e^{-\int_t^Tr_s\,ds}X] = P(t, T)\,\E^{T}_t[X]$. Le prix forward
$V_t/P(t,T)$ de tout actif négocié est une $\mathbb Q^T$-martingale, et en particulier $\E^T[r_T] = f(0, T)$,
le taux forward instantané.
\end{proposition}

\begin{proof}
Appliquer la formule de valorisation avec $\mathcal N = P(\cdot, T)$, pour laquelle $P(T, T) = 1$. Pour la dernière affirmation,
$f(0,T) = -\partial_T\ln P(0,T) = \E^{\mathbb Q}[r_Te^{-\int_0^Tr}]/P(0,T) = \E^T[r_T]$.
\end{proof}

L'espérance du produit d'un facteur d'actualisation stochastique et d'un gain est devenue un
facteur d'actualisation déterministe multiplié par une espérance. Pour le taux court d'Ornstein--Uhlenbeck du chapitre~4
($r_0 = 3\%$, $\bar r = 4\%$, $\kappa = 0.5$, $\sigma = 1\%$), l'obligation a pour volatilité $-\sigma B(T - t)$, donc par
le \cref{thm:qm:girsanov-and-changes-of-numeraire:change}, la dérive du taux court sous $\mathbb Q^T$ perd
$\sigma^2B(T - t)$~: $r_T$ reste gaussien, avec une moyenne inférieure. $\E^{\mathbb Q}[r_5] = 3.918\%$ tandis que $\E^5[r_5] = f(0,
5) = 3.901\%$. L'écart, effet de convexité, croît jusqu'à $\sigma^2/2\kappa^2 = 2$ points de base aux longues échéances
(la \cref{fig:qm:girsanov-and-changes-of-numeraire:convexity}).

\begin{omfigure}
\begin{tikzpicture}
\begin{axis}[width=11cm,height=4.6cm,
  xlabel={échéance $t$ (années)},ylabel={$\E^{\mathbb Q}[r_t] - f(0,t)$ (pb)},
  tick label style={font=\small},label style={font=\small},
  xmin=0,xmax=30,ymin=0,ymax=2.2,grid=major,grid style={gray!25},
  table/col sep=comma]
\addplot[omDef,very thick] table[x=t,y=gap]{figdata/methods/05-girsanov-and-changes-of-numeraire/convexity.csv};
\addplot[gray,dashed,domain=0:30] {2};
 \node[font=\footnotesize,anchor=north east,fill=white,inner sep=1pt] at (axis cs:29,1.9) {$\sigma^2/2\kappa^2 = 2$~pb};
\end{axis}
\end{tikzpicture}

\omcaption{Taux court espéré sous la \omterm{def:qm:girsanov-and-changes-of-numeraire:emm}{probabilité risque-neutre} moins le taux forward instantané,
qui est le taux court espéré sous la \omterm{def:qm:girsanov-and-changes-of-numeraire:forward}{probabilité forward} de même échéance, pour le
taux court d'Ornstein--Uhlenbeck du chapitre. L'écart, $\sigma^2B(t)^2/2$, est le prix de la
convexité de l'obligation. Données~: formules fermées, tutoriel du chapitre.}\label{fig:qm:girsanov-and-changes-of-numeraire:convexity}
\end{omfigure}

Le caplet de l'exemple d'ouverture verse $\delta(F - K)^+$ en $T_1 + \delta$ sur le taux $F$ fixé en $T_1 = 5$ pour $\delta = \tfrac14$
an, avec $K = 4.5\%$. Sa valeur en $T_1$ est $(1 + \delta K)(x - P(T_1, T_1 + \delta))^+$ avec $x = 1/(1 + \delta K) =
0.98888$~: un put sur une obligation zéro-coupon. Sous $\mathbb Q^{T_1}$, $P(T_1, T_1 + \delta) = e^{A - Br_{T_1}}$ est log-normale,
et une formule de type Black en découle (Jamshidian, 1989) avec une volatilité de l'obligation $\sigma_P =
\sigma\sqrt{(1 - e^{-2\kappa T_1})/2\kappa}\,B(\delta) = 0.234\%$~: 3.26 points de base de notionnel, soit 32\,600~dollars sur
100~millions de dollars. La \Cref{fig:qm:girsanov-and-changes-of-numeraire:mc} compare les deux simulations.

\begin{omfigure}
\begin{tikzpicture}
\begin{axis}[width=11cm,height=5cm,xmode=log,log basis x=2,
  xlabel={trajectoires},ylabel={caplet (pb du notionnel)},
  tick label style={font=\small},label style={font=\small},
  xmin=180,xmax=90000,ymin=1.8,ymax=4.4,grid=major,grid style={gray!25},
  legend columns=3,legend style={font=\footnotesize,at={(0.5,-0.3)},anchor=north,draw=none,fill=none,/tikz/every even column/.append style={column sep=8pt}},
  table/col sep=comma]
\addplot[omDef,thick,mark=*,mark size=1.5pt,error bars/.cd,y dir=both,y explicit]
  table[x=n,y=q,y error plus expr=\thisrow{qhi}-\thisrow{q},y error minus expr=\thisrow{q}-\thisrow{qlo}]{figdata/methods/05-girsanov-and-changes-of-numeraire/convergence.csv};
\addplot[omThm,thick,mark=triangle*,mark size=1.8pt,error bars/.cd,y dir=both,y explicit]
  table[x expr=\thisrow{n}*1.12,y=fwd,y error plus expr=\thisrow{fhi}-\thisrow{fwd},y error minus expr=\thisrow{fwd}-\thisrow{flo}]{figdata/methods/05-girsanov-and-changes-of-numeraire/convergence.csv};
\addplot[black,dashed,domain=180:90000] {3.2586};
\legend{sous $\mathbb Q$ (260 normales par trajectoire),sous $\mathbb Q^{T_1}$ (une normale),formule fermée}
\end{axis}
\end{tikzpicture}

\omcaption{Le caplet par Monte-Carlo sous la \omterm{def:qm:girsanov-and-changes-of-numeraire:emm}{probabilité risque-neutre} (pas hebdomadaires, actualisation au
\omterm{def:qm:girsanov-and-changes-of-numeraire:numeraire}{compte monétaire}) et sous la \omterm{def:qm:girsanov-and-changes-of-numeraire:forward}{probabilité forward} à la date de fixation, avec $\pm 2$ erreurs types, face à la
formule fermée de 3.26 pb. À nombre égal de trajectoires, les erreurs sont les mêmes~; la trajectoire sous \omterm{def:qm:girsanov-and-changes-of-numeraire:forward}{probabilité forward} coûte
une normale au lieu de 260. Données~: le tutoriel du chapitre, à graine fixée.}\label{fig:qm:girsanov-and-changes-of-numeraire:mc}
\end{omfigure}

Le gain ne tient pas à une variance plus faible par trajectoire~: le facteur d'actualisation varie peu
au regard du payoff, et les deux erreurs types par trajectoire ne diffèrent l'une de l'autre que de moins d'un demi-pour-cent. Il tient
à la dimension du problème, qui passe de 260 normales à une seule, puis à l'absence totale de simulation.

\section{Probabilité annuité}

\begin{definition}[Probabilité annuité]\label{def:qm:girsanov-and-changes-of-numeraire:annuity}
Pour un swap à paiements fixes aux dates $T_1 < \dots < T_n$ et de fractions d'accumulation $\delta_i$, la \emph{probabilité
annuité}\index{probabilité annuité} $\mathbb Q^A$ est la \omterm{def:qm:girsanov-and-changes-of-numeraire:emm}{probabilité martingale équivalente} pour le \omterm{def:qm:girsanov-and-changes-of-numeraire:numeraire}{numéraire} $A_t =
\sum_i\delta_iP(t, T_i)$, l'annuité du livre~2, chapitre~9.
\end{definition}

Le taux de swap au pair est $S_t = (P(t, T_0) - P(t, T_n))/A_t$, une valeur de marché divisée par le \omterm{def:qm:girsanov-and-changes-of-numeraire:numeraire}{numéraire}~:
une $\mathbb Q^A$-martingale. Une swaption payeuse exercée en $T_0$ verse $A_{T_0}(S_{T_0} - K)^+$, donc sa valeur est $A_0\,\E^A[(S_{T_0}
- K)^+]$~: avec une loi normale pour $S_{T_0}$, il s'agit de la formule de Bachelier du livre~2,
chapitre~13, et avec une loi lognormale, celle de Black. Le \omterm{def:qm:girsanov-and-changes-of-numeraire:changenum}{changement de numéraire} remplace une swaption, un droit portant sur
toute une courbe, par un call sur une seule \omterm{def:qm:probability-at-speed:martingale}{martingale}. Le livre~6 construit ses modèles de taux sur cette
probabilité et sur la \omterm{def:qm:girsanov-and-changes-of-numeraire:forward}{probabilité forward}.

\section{Tutoriel~: deux probabilités, un seul prix}\label{tut:qm:girsanov-and-changes-of-numeraire:caplet}

\begin{tutorial}
\textbf{Objectif.} Évaluer le caplet sous deux probabilités, vérifier la repondération entre elles, et voir
Girsanov à l'œuvre sur un histogramme. \textbf{État final~:}
\cref{fig:qm:girsanov-and-changes-of-numeraire:girsanov,fig:qm:girsanov-and-changes-of-numeraire:convexity,fig:qm:girsanov-and-changes-of-numeraire:mc}
et le prix de 3.26~pb.
\begin{tutsteps}
\item \textbf{Les poids.} Le projet en cours convertit les espérances entre \omterm{def:qm:girsanov-and-changes-of-numeraire:numeraire}{numéraires} et construit
des densités de Girsanov.
\omcode{code/firm/numeraire/firm_numeraire.py}{21}{37}{Poids du changement de numéraire, moyenne pondérée et densité de Girsanov.}\label{lst:qm:girsanov-and-changes-of-numeraire:weights}
\item \textbf{Le générateur sous la \omterm{def:qm:girsanov-and-changes-of-numeraire:forward}{probabilité forward}}~: la moyenne de $r_{T}$ perd $\sigma^2\int_0^Te^{-\kappa(T-s)}B(T - s)\,ds$.
\omcode{code/methods/05-girsanov-and-changes-of-numeraire/python/qm_numeraire.py}{85}{99}{Le taux court à la date de fixing sous la probabilité forward, et le caplet par une normale par trajectoire.}\label{lst:qm:girsanov-and-changes-of-numeraire:forward}
\item \textbf{Exécuter} \texttt{problem()}, \texttt{reweighting\_check()} (100\,000 trajectoires risque-neutres
repondérées par $d\mathbb Q^{T}/d\mathbb Q$ donnent $\E^{T}[r_5] = 3.906\% \pm 0.003\%$ contre $3.901\%$) et
\texttt{fig\_numeraire.py}.
\end{tutsteps}
\textbf{Pour aller plus loin.} Évaluer le caplet sous la probabilité de la date de paiement $\mathbb Q^{T_1 + \delta}$, sous
laquelle le taux forward lui-même est une \omterm{def:qm:probability-at-speed:martingale}{martingale}~; évaluer une swaption sous la \omterm{def:qm:girsanov-and-changes-of-numeraire:annuity}{probabilité annuité} et
comparer avec une simulation risque-neutre de la courbe entière.
\end{tutorial}

\section{Construction~: passer d'un numéraire à l'autre}\label{bld:qm:girsanov-and-changes-of-numeraire:numeraire}

\begin{build}
\bfield{Rôle} Chaque moteur de valorisation de la firme miniature énonce son \omterm{def:qm:girsanov-and-changes-of-numeraire:numeraire}{numéraire}, et chaque
prix simulé est confronté à lui~: un prix divisé par le \omterm{def:qm:girsanov-and-changes-of-numeraire:numeraire}{numéraire} doit être une \omterm{def:qm:probability-at-speed:martingale}{martingale}.
\bfield{Interface} \texttt{change\_of\_numeraire\_weights(m\_T, m\_0, n\_T, n\_0)}~;
\texttt{reweighted\_mean(x, w)}~; \texttt{girsanov\_weights(dw, gamma, dt)}~;
\texttt{deflated\_martingale\_test(prices, numeraire)} renvoie le plus grand $|t|$ sur les dates.
\bfield{Règles} Les poids ont une espérance égale à un sous la probabilité d'échantillonnage (un test, non une hypothèse)~;
l'intégrande d'une densité de Girsanov est adaptée~; le test de \omterm{def:qm:probability-at-speed:martingale}{martingale} utilise des trajectoires simulées sous
la probabilité propre au \omterm{def:qm:girsanov-and-changes-of-numeraire:numeraire}{numéraire}.
\bfield{Tests d'acceptation} \texttt{code/firm/numeraire/tests/}~: les poids de Cameron--Martin déplacent la
moyenne d'une loi gaussienne et non sa variance~; la repondération des trajectoires de taux court risque-neutres reproduit la
moyenne sous la \omterm{def:qm:girsanov-and-changes-of-numeraire:forward}{probabilité forward}~; les prix actualisés des obligations passent le test de \omterm{def:qm:probability-at-speed:martingale}{martingale}, les prix non actualisés échouent.
\bfield{Pour aller plus loin} La probabilité spot d'un compte renouvelé à pas discrets (livre~6)~; les sensibilités par ratio de vraisemblance,
qui sont des poids de Girsanov dérivés par rapport à un paramètre (livre~5,
chapitre~23).
\end{build}

\begin{omsources}
\item I.~V.~Girsanov, ``On transforming a certain class of stochastic processes by absolutely
continuous substitution of measures'', \emph{Theory of Probability and its Applications} 5, 1960.
\item R.~H.~Cameron and W.~T.~Martin, ``Transformations of Wiener integrals under translations'',
\emph{Annals of Mathematics} 45, 1944.
\item H.~Geman, N.~El Karoui and J.-C.~Rochet, ``Changes of numéraire, changes of probability
measure and option pricing'', \emph{Journal of Applied Probability} 32, 1995.
\item F.~Jamshidian, ``An exact bond option formula'', \emph{Journal of Finance} 44, 1989.
\item W.~Margrabe, ``The value of an option to exchange one asset for another'', \emph{Journal of
Finance} 33, 1978.
\end{omsources}

\section{Exercices}

\begin{exercise}[$\star$]\label{exo:qm:girsanov-and-changes-of-numeraire:1}
Sous $d\mathbb Q/d\P = \exp(\gamma W_T - \gamma^2T/2)$ avec $\gamma = 0.5$ et $T = 4$, quelle est la loi de $W_T$~?
\end{exercise}

\begin{exercise}[$\star$]\label{exo:qm:girsanov-and-changes-of-numeraire:2}
Une action a un rendement espéré de 9~\% et une volatilité de 25~\%, et le taux court vaut 4~\%. Calculer sa dérive sous
la \omterm{def:qm:girsanov-and-changes-of-numeraire:emm}{probabilité risque-neutre}, et le $\gamma$ qu'utilise Girsanov.
\end{exercise}

\begin{exercise}[$\star$]\label{exo:qm:girsanov-and-changes-of-numeraire:3}
Un payoff à cinq ans a $\E^{5}[X] = 2.4$. Avec le $P(0, 5)$ du chapitre, calculer sa valeur aujourd'hui.
\end{exercise}

\begin{exercise}[$\star\star$]\label{exo:qm:girsanov-and-changes-of-numeraire:4}
Expliquer pourquoi $\E^{\mathbb Q}[r_T] > f(0, T)$ dans le modèle du chapitre, et calculer l'écart en $T = 5$ à partir de
$\sigma^2B(T)^2/2$.
\end{exercise}

\begin{exercise}[$\star\star$]\label{exo:qm:girsanov-and-changes-of-numeraire:5}
Montrer que le taux de swap au pair est une \omterm{def:qm:probability-at-speed:martingale}{martingale} sous la \omterm{def:qm:girsanov-and-changes-of-numeraire:annuity}{probabilité annuité}, et écrire la valeur d'une
swaption receveuse comme une espérance sous cette probabilité.
\end{exercise}

\begin{exercise}[$\star\star$]\label{exo:qm:girsanov-and-changes-of-numeraire:6}
Deux actions suivent des \omterm{def:qm:stochastic-differential-equations:gbm}{mouvements browniens géométriques} de volatilités 20~\% et de corrélation 0.5, toutes deux valant 100,
sans dividende. Avec la seconde action comme \omterm{def:qm:girsanov-and-changes-of-numeraire:numeraire}{numéraire}, évaluer l'option de recevoir la première action en échange de
la seconde action dans un an.
\end{exercise}

\begin{exercise}[$\star\star\star$]\label{exo:qm:girsanov-and-changes-of-numeraire:7}
Déduire la probabilité de premier passage avec dérive de la \cref{prop:qm:brownian-motion:driftpassage}~: éliminer
la dérive par Girsanov, appliquer la loi conjointe du minimum et du point final fournie par le principe de réflexion,
puis intégrer. Le vérifier sur l'arrêt du chapitre~2 avec la dérive $-\sigma^2/2$.
\end{exercise}

\begin{exercise}[$\star\star\star$]\label{exo:qm:girsanov-and-changes-of-numeraire:8}
\emph{Trouver la faille.} «~Pour valoriser l'option, nous simulons l'action avec la dérive
de 12~\% par an prévue par notre équipe de recherche, et nous actualisons les payoffs au taux sans risque~: la prévision est
notre edge.~»
\end{exercise}

\section{Problème~: deux probabilités, un seul prix}

\begin{problem}[{Problème du week-end --- un caplet à cinq ans, de deux façons}]\label{pb:qm:girsanov-and-changes-of-numeraire:1}
Le taux court suit $dr = \kappa(\bar r - r)\,dt + \sigma\,dW^{\mathbb Q}$ avec $r_0 = 3\%$, $\bar r = 4\%$, $\kappa = 0.5$ et
$\sigma = 1\%$. Un caplet fixe en $T_1 = 5$ sur le taux de $[5, 5.25]$, paie en 5.25, strike 4.5~\%,
notionnel 100 millions de dollars.

\medskip\noindent\textbf{Partie I --- La courbe.}
\begin{enumerate}
\item Quels sont $P(0, 5)$ et $P(0, 5.25)$~?
\item Quel est le taux forward sur $[5, 5.25]$~?
\item Quels sont $\E^{\mathbb Q}[r_5]$ et $f(0, 5)$, et pourquoi diffèrent-ils~?
\item Écrire le caplet comme un put sur une obligation zéro-coupon~: strike et nombre d'obligations.
\item Quelle est la volatilité de l'obligation $\sigma_P$ dans la formule de Jamshidian~?
\end{enumerate}

\medskip\noindent\textbf{Partie II --- Trois prix.}
\begin{enumerate}[resume]
\item Que donne la forme fermée, en points de base et en dollars~?
\item Combien de normales utilise un chemin risque-neutre avec des pas hebdomadaires, et un chemin en \omterm{def:qm:girsanov-and-changes-of-numeraire:forward}{probabilité forward}~?
\item Que donnent 20\,000 chemins risque-neutres, avec leur erreur type~?
\item Et 20\,000 chemins en \omterm{def:qm:girsanov-and-changes-of-numeraire:forward}{probabilité forward}~?
\item Les deux sont-ils dans leur erreur type de la forme fermée~?
\end{enumerate}

\medskip\noindent\textbf{Partie III --- Ce que le \omterm{def:qm:probability-at-speed:change}{changement de probabilité} a apporté.}
\begin{enumerate}[resume]
\item Quelle est la densité $d\mathbb Q^{T_1}/d\mathbb Q$ sur un chemin~?
\item Que donne la repondération de 100\,000 chemins risque-neutres pour $\E^{T_1}[r_5]$~?
\item Quel est le rapport des deux erreurs types à nombre de chemins égal, et pourquoi est-il proche de un~?
\item Quel est alors le gain à précision égale~?
\item Avec des pas quatre fois plus fins sous $\mathbb Q$, le prix risque-neutre change-t-il~?
\end{enumerate}

\medskip\noindent\textbf{Partie IV --- Jugement.}
\begin{enumerate}[resume]
\item Un cap est une suite de caplets aux dates de paiement différentes. Quelle probabilité utilisez-vous~?
\item Quelle probabilité price une swaption, et qu'est-ce qui devient une \omterm{def:qm:probability-at-speed:martingale}{martingale}~?
\item Que doit vérifier un test d'acceptation d'une bibliothèque de pricing concernant son \omterm{def:qm:girsanov-and-changes-of-numeraire:numeraire}{numéraire}~?
\item Énoncer le \emph{résultat nommé}~: le prix du caplet et ce que la \omterm{def:qm:girsanov-and-changes-of-numeraire:forward}{probabilité forward} a apporté.
\item En une phrase~: que change un \omterm{def:qm:girsanov-and-changes-of-numeraire:changenum}{changement de numéraire}, et que ne change-t-il pas~?
\end{enumerate}
\end{problem}

\section{Questions d'entretien}

\begin{interviewq}[$\star$ \iqroles{chercheur, banque}]\label{iq:qm:girsanov-and-changes-of-numeraire:1}
Que change le théorème de Girsanov à un \omterm{def:qm:brownian-motion:bm}{mouvement brownien}, et que ne peut-il pas changer~?
\end{interviewq}

\begin{interviewq}[$\star$ \iqroles{banque, trader}]\label{iq:qm:girsanov-and-changes-of-numeraire:2}
Pourquoi une action dérive-t-elle au taux sans risque sous la \omterm{def:qm:girsanov-and-changes-of-numeraire:emm}{probabilité risque-neutre}, alors que personne ne s'y attend~?
\end{interviewq}

\begin{interviewq}[$\star\star$ \iqroles{banque}]\label{iq:qm:girsanov-and-changes-of-numeraire:3}
Qu'est-ce que la \omterm{def:qm:girsanov-and-changes-of-numeraire:forward}{probabilité forward}, et pourquoi le taux forward est-il une \omterm{def:qm:probability-at-speed:martingale}{martingale} sous celle-ci~?
\end{interviewq}

\begin{interviewq}[$\star\star$ \iqroles{banque, chercheur}]\label{iq:qm:girsanov-and-changes-of-numeraire:4}
Pricez l'option d'échanger l'actif 2 contre l'actif 1 en $T$. Quel \omterm{def:qm:girsanov-and-changes-of-numeraire:numeraire}{numéraire} choisissez-vous, et pourquoi~?
\end{interviewq}

\begin{interviewq}[$\star\star$ \iqroles{développeur, banque}]\label{iq:qm:girsanov-and-changes-of-numeraire:5}
Comment testeriez-vous qu'un moteur de pricing Monte Carlo utilise une probabilité cohérente~?
\end{interviewq}

\begin{interviewq}[$\star\star\star$ \iqroles{chercheur}]\label{iq:qm:girsanov-and-changes-of-numeraire:6}
Estimez $\P(W_1 > 5)$ par simulation avec moins de 10\,000 tirages.
\end{interviewq}
